\section{De la T2V básica a la edición, la personalización y el audio}

Las regularidades visuales que aprende el modelo base pueden sustentar diversas aplicaciones, pero cada tarea requiere una interfaz de entrada y unos datos de entrenamiento distintos. Esta sección distingue seis estados representativos: la generación a partir de solo texto se centra en la plausibility espaciotemporal; la edición, en la conservación del contenido; la personalización, en la identidad; el audio, en la sincronización entre modalidades. Llamarlos a todos «Movie Gen» puede ocultar las diferencias de entrenamiento dentro de la familia de modelos.

\subsection{¿El modelo aprendió realmente el «mundo visual»?}

Las secciones anteriores sobre arquitectura y datos explicaron cómo se entrena el modelo; esta sección vuelve a las salidas para examinar qué aprendió en realidad. El ponente usa dos ejemplos para mostrar que un modelo a gran escala puede generar movimientos de cámara, movimientos de objetos y ciertas relaciones físicas. Al leer estas figuras conviene buscar continuidad e interacciones, no solo la textura de un fotograma aislado; al mismo tiempo, hay que respetar los límites de la evidencia: los datos observacionales pueden producir un ajuste sólido de regularidades, pero no generan automáticamente un modelo del mundo causal sobre el que se pueda intervenir.

\lecturefigure{slide-257-learned-visual-world.jpg}{Ejemplo de la carrera con papalote: el sujeto, el papalote, la playa y el movimiento de cámara evolucionan en conjunto}{00:41:17}

Este fragmento exige que el modelo mantenga la identidad y la ropa de la persona, que el papalote se coordine con la dirección del movimiento y que resuelva el seguimiento de la cámara. Muestra una generación conjunta de múltiples factores, pero no permite saber si variables latentes como la fuerza del viento o la tensión del hilo están representadas de forma explícita.

\lecturefigure{slide-305-physics-sample.jpg}{Ejemplo de la alberca: la superficie del agua, la dona, el perezoso y la bebida forman relaciones de contacto y de flotación}{00:42:11}

Que las relaciones de contacto sean visualmente razonables indica que el modelo aprendió una gran cantidad de sentido común estadístico. Una evaluación física más rigurosa debería poner a prueba contrafactuales: cambiar la masa de los objetos, el líquido, la velocidad de colisión o las restricciones, y observar si el resultado sigue las leyes de manera estable, en lugar de solo parecer razonable en escenas comunes.

\begin{warningbox}{«World model» tiene al menos tres grados de exigencia}
La definición débil es poder generar futuros visualmente razonables; la intermedia exige mantener la coherencia ante intervenciones y contrafactuales; la fuerte exige además que sirva para planificar y predecir. Los ejemplos de Movie Gen respaldan el primer nivel y rozan fenómenos del segundo, pero no basta con demostraciones para afirmar que se cuenta con un simulador causal completo.
\end{warningbox}

\teachervoice{00:41:11--00:42:52, el docente dice que el modelo parece haber aprendido objetos, movimiento y leyes físicas a partir de ver videos; lo considera una sorpresa del escalamiento, no una demostración de razonamiento físico explícito.}

\subsection{Edición precisa: además de la calidad de generación, hay que preservar el video original}

La generación a partir de solo texto solo necesita construir un mundo razonable; la edición, además, debe respetar un mundo que ya existe. Por eso esta sección divide la evaluación en «qué se cambió correctamente» y «qué se cambió por error»: el modelo de edición recibe a la vez el video original y la instrucción de texto, y debe limitar los cambios a la región y los atributos objetivo. Las dos cuadrículas de cuatro paneles que siguen muestran sustituciones locales, eliminaciones, adiciones y cambios de estilo; son más adecuadas que una salida única para comprobar qué contenido se conserva.

\lecturefigure{slide-337-precise-editing-vr.jpg}{Edición de una escena de VR: sustituir, retirar el visor y añadir efectos visuales}{00:42:54}

Que una misma entrada admita múltiples instrucciones indica que la interfaz de edición puede desacoplar el source content del edit intent. La evaluación debe separarse en edit fidelity y preservation: la primera examina el cambio objetivo; la segunda, si la persona, la acción, el fondo y la toma derivan de forma no intencionada.

\lecturefigure{slide-374-precise-editing-penguin.jpg}{Edición del video de pingüinos: vestimenta, objetos del entorno y estilo de dibujo a lápiz}{00:43:54}

Las ediciones de estilo global y las de objetos locales tienen dificultades distintas. Añadir un paraguas es una modificación estructural local; ponerles ropa exige sincronizarla con el movimiento; el estilo a lápiz vuelve a dibujar todo el cuadro, pero debe conservar la disposición de la escena. Una puntuación única ocultaría estos modos de falla.

\begin{knowledgebox}{Por qué escasean los datos de edición}
En el mundo real casi no existen pairs naturales de «el mismo video antes/después de la edición + instrucción precisa». En clase se indicó que Movie Gen usa una construcción autosupervisada para producir la señal de entrenamiento. Sin este diseño de datos, un modelo T2V básico, por bueno que sea generando, difícilmente conservaría con precisión el contenido de entrada.
\end{knowledgebox}

\subsection{La personalización y el audio son rutas de condicionamiento distintas}

Después de la tarea de edición, esta sección compara dos extensiones de condicionamiento completamente diferentes: el video personalizado añade una reference identity y la generación de audio añade un condicionamiento video/text-to-audio. Ambas amplían la capacidad de generación básica, pero persiguen objetivos distintos: la personalización debe equilibrar identidad y acción; el audio debe alinearse con los eventos visuales, el ritmo y la semántica del entorno. Por eso no pueden resumirse con una misma puntuación de «calidad de video».

\lecturefigure{slide-388-personalization.jpg}{Video personalizado: el retrato de referencia y el prompt de texto determinan juntos el sujeto y la escena}{00:44:12}

La imagen de referencia de la esquina superior derecha aporta la información de identidad, y el video principal debe conservar esa identidad con nuevas poses, iluminación y escenas. Copiar sin más la textura del rostro produce un movimiento rígido; obedecer en exceso al texto, en cambio, hace perder la identidad. Hay que medir a la vez identity similarity, motion quality y prompt adherence.

\lecturefigure{slide-419-synchronized-audio.jpg}{Movie Gen Audio genera truenos y música de fondo sincronizados a partir del video y el texto}{00:44:43}

El modelo de audio necesita saber cuándo cae el rayo, cuál es el ánimo de la escena y qué sonidos pertenecen al interior del cuadro (diegetic) o al exterior. La conferencia señala explícitamente que el modelo T2V mostrado antes produce por sí mismo video sin sonido; el audio sincronizado proviene de un modelo independiente, Movie Gen Audio.

\begin{importantbox}{Movie Gen es una familia de modelos, no una sola red}
La T2V básica, video editing, personalization y audio generation comparten una línea de investigación, pero difieren en entradas, objetivos, datos y posentrenamiento. Al evaluar o desplegar, hay que indicar qué modelo y qué pipeline se usan; no se pueden atribuir todas las capacidades de la familia al backbone T2V de 30B.
\end{importantbox}

\teachervoice{01:11:28--01:13:20, al final el docente vuelve a aclarar: el modelo principal de la clase genera video sin sonido y el audio proviene de un modelo independiente. En teoría, la generación conjunta de audio y video podría compartir información, pero los datos de alineación de alta calidad son más escasos que los de video por sí solo.}

\subsection{Resumen de la sección}

La familia Movie Gen abarca la generación pura, la edición, la personalización y el audio, pero estas capacidades no provienen de un mismo objetivo de entrenamiento. La T2V básica muestra un ajuste sólido de regularidades visuales; la edición enfatiza la conservación del contenido; la personalización, la identidad; el audio, la sincronización entre modalidades. El nombre de cada capacidad debe vincularse al modelo y a la ruta de datos concretos.
